\chapter{Нелинейная проверка отрицательной моды и решёточное закрепление вихря}

\section{Почему гессиана sampled-фона недостаточно}

Предыдущая глава получила отрицательное значение $-0.189047\ldots$ на
профиле, который был решён как непрерывная радиальная краевая задача, а
затем только выбран на декартовой сетке. Такой sampled-профиль не является
точной стационарной точкой дискретной энергии: её градиент мал, но не
равен нулю. Гессиан в нестационарной точке не даёт сертификата
устойчивости.

Дополнительная опасность состоит в нарушении непрерывной трансляционной
симметрии сеткой. В дискретных системах центрированный в узле и
центрированный между узлами дефекты разделяет барьер Пайерлса--Набарро.
Различие таких ветвей и исчезновение барьера в трансляционно согласованных
дискретизациях обсуждаются в \path{arXiv:1601.04598} и
\path{arXiv:nlin/0603047}; вихревое проявление решёточного потенциала
описано в \path{arXiv:0902.1201}.

\section{Одномерный нелинейный срез}

Вдоль отрицательного собственного вектора сетки $33$ полная энергия с
положительным калибровочным квадратом образует симметричный срез. Минимум
положительной ветви достигается при метрической амплитуде
\begin{equation}
 q_*=0.3207811\ldots,
\end{equation}
а выигрыш равен
\begin{equation}
 \Delta E(q_*)=-0.0129845\ldots.
\end{equation}
Однако полная релаксация этой ветви перемещает распределение потока
примерно на один шаг сетки. Это характерно для перехода между
узловой и межузловой реализациями одного дефекта, а не для рождения новой
континуальной конфигурации.

\section{Сначала стационарная точка, затем гессиан}

Для строгой проверки использованы нечётные сетки $25,33,41$, содержащие
центральный узел. В нём закреплено условие
\begin{equation}
 \phi(0)=0,
\end{equation}
которое удаляет перемещение ядра, но сохраняет все внутренние
возмущения. На каждой сетке сначала минимизирована полная дискретная
энергия вместе с неотрицательным калибровочным квадратом, и только после
сходимости построен новый гессиан.

Среднеквадратичные нормы проецированного градиента не превосходят
$4.7\cdot10^{-9}$. Нижние собственные значения релаксированного оператора
равны
\begin{equation}
\begin{array}{c|ccc}
 N & 25 & 33 & 41\\ \hline
 \lambda_{\min} &
 0.01202263 & 0.00086613 & 0.00047270
\end{array}
\end{equation}
и отрицательных мод среди шести проверенных нет ни на одной сетке.
Уменьшение нижней пары к нулю согласуется с восстановлением переносной
моды в континуальном пределе.

Центр потока релаксированного решения смещён от закреплённого ядра на
величину порядка шага сетки. Его координата уменьшается вместе с шагом:
примерно $0.349$, $0.264$, $0.197$ при $h=0.833$, $0.625$, $0.5$.
Следовательно, смещение не задаёт конечного физического масштаба.

\section{Исправленный статус}

Отрицательная мода предыдущей главы воспроизводима как свойство
нестационарного sampled-фона, но не переживает правильный порядок
операций
\begin{equation}
 \text{дискретная стационаризация}
 \longrightarrow
 \text{фиксация переносов}
 \longrightarrow
 \text{гессиан}.
\end{equation}
Поэтому она не является доказанной физической неустойчивостью прямой
нити и не насыщается новым конечным дефектом.

\begin{nogocriterion}
Нельзя объявлять вихрь полностью устойчивым: проверена только эффективная
пятикомпонентная подсистема, а нижние положительные уровни ещё не
исследованы в трансляционно ковариантной дискретизации. Тем более нельзя
переходить к хопфовой петле или рождению материи без полного оператора
$Q+T+B$.
\end{nogocriterion}

\begin{protocol}
\begin{enumerate}
 \item одномерная двойная яма: \textbf{воспроизведена};
 \item её нелинейный исход: \textbf{сдвиг центра по сетке};
 \item центрированные стационарные решения: \textbf{получены};
 \item отрицательные моды после релаксации: \textbf{не найдены};
 \item прежний отрицательный сертификат: \textbf{снят};
 \item полная устойчивость $Q+T+B$: \textbf{не закрыта};
 \item следующий гейт: \textbf{трансляционно ковариантная дискретизация}.
\end{enumerate}
\end{protocol}

Машинный сертификат сохранён в
\path{s2t_v6_bosonic_defect_corrected_vortex_negative_mode_nonlinear_saturation_gate_results.json}.